\lecture{3}
\title{Regular Expressions, The Pumping Lemma}

\begin{document}

\subsection{Generalized NFAs}

Recall that we proved in
\href{/18.404/lectures/2/#regular-expressions-are-regular-languages}{Lecture 2}
the following theorem
by introducing a new kind of FA (the NFA).

\begin{theorem}{Regularness}
    Every regular expression is a regular language.
\end{theorem}

In this lecture, we'll introduce \emph{another} kind of FA---the
\textbf{generalized NFA (GNFA)}---to prove the reverse direction.

\begin{theorem}{Regularness}
    Every regular language is a regular expression.
\end{theorem}

\textbf{Definition.}
A generalized NFA is an NFA in which edges may be labeled
with any regular expression (not just atoms!).
We say a GNFA is \textbf{good} if it satisfies the following properties:
\begin{itemize}
    \item There is a single accept state distinct from the start state.
    \item There exists a transition between every pair of states
    and a self-loop at every state (possibly with label $\emptyset$) \dots
    \begin{itemize}
        \item \dots except the start state must have indegree zero \dots
        \item \dots and the accept state must have outdegree zero.
    \end{itemize}
\end{itemize}
Evidently, every GNFA can be easily translated into
a good GNFA by adding $\emptyset$ edges and some additional states.
\begin{center}
    \frame{\includegraphics[width=0.75\textwidth]{lecture-03/GNFA-examples}}
\end{center}

Expectedly, good GNFAs are equally as representative of regular languages
as ordinary NFAs (or DFAs) are.

\begin{theorem}{GNFA Implies NFA}
    A language $A$ is recognized by some NFA
    if and only if $A$ is recognized by some GNFA.
\end{theorem}

\begin{proof}
    The NFA $\implies$ GNFA direction is trivial,
    so we only show the GNFA $\implies$ NFA direction.

    Simply apply the fact from \href{/18.404/lectures/2}{Lecture 2}
    that every regular expression is the language of some NFA.
    That is, for each edge with regular expression $r$,
    replace that edge with the NFA whose language is $r$,
    linking into and out of the NFA using $\epsilon$-edges.
    \begin{center}
        \frame{\includegraphics[width=0.5\textwidth]{lecture-03/GNFA-equivalence}}
    \end{center}
    If this new NFA accepts a string $s$,
    then the original GNFA must also accept $s$;
    to see why, consider breaking up the computation path of $s$
    into sub-paths $p \to q$ where $p$ and $q$ are states of the original GNFA,
    and so\dots ~ $\blacksquare$
\end{proof}

\subsection{Regular Languages are Regular Expressions}

We now prove the following theorem.

\begin{theorem}{GNFA Regularness}
    Every good GNFA regular language is a regular expression.
\end{theorem}

\begin{proof}
    Proof by induction on the number of states.
    The base case is trivial.

    For the inductive step, remove a state $X$ from the GNFA $G$
    that is not the start or accept state.
    \begin{center}
        \frame{\includegraphics[width=0.5\textwidth]{lecture-03/GNFA-proof}}
    \end{center}
    Furthermore, for \emph{every} ordered pair of states
    $(p, q) \in (Q \setminus \{X\})^2$,
    update the label of the edge $p \to q$ as shown above.
    This leaves the language of $G$ unchanged,
    so we can induct down successfully. ~ $\blacksquare$
\end{proof}

\begin{remark}
    CYU: Where, in the proof, does the ``goodness'' of the GNFA get used?
\end{remark}

And of course, by combining this theorem with the
``GNFA Implies NFA'' theorem preceding it, we get:

\begin{theorem}{DFA Regularness}
    Every (DFA) regular language is a regular expression.
\end{theorem}

\subsection{Proving Nonregularity}

All of our work so far has gone toward proving
the regularity of languages.
What about nonregularity?

\textbf{Example.}
Consider the following languages with alphabet $\Sigma = \{\str{0}, \str{1}\}$.
\[
B := \{w \mid w
\text{ has as many } \str{0}\text{s} \text{ as } \str{1}\text{s}\}
~~~ \parallel ~~~
C := \{w \mid w
\text{ has as many } \str{01}\text{s} \text{ as } \str{10}\text{s}\}
\]
As we'll see, $B$ is nonregular.
But $C$ is regular---consider
a five-state FA that tracks
$(\textsc{First Symbol}, \textsc{Last Symbol})$.

\begin{theorem}{Pumping Lemma}
    For every regular language $A$,
    there must be some \textbf{pumping length} $p$
    such that:
    \begin{quote}
        For all $s \in A$ with $|s| \geq p$,
        we may write $s = xyz$ with $y \neq \epsilon$
        such that $|xy| \leq p$
        and $xy^iz \in A$ for all $i \in \mathbb{N}$.
    \end{quote}
\end{theorem}

\begin{proof}
    The pumping length of $A = L(M)$ is the number of states in the FA $M$.

    Consider some string $s \in A$ with $|s| \geq p$.
    Suppose it begins at $S$ and ends at $T$ through some path $\mathcal{P}$.
    Since $|s| \geq p$, the path $\mathcal{P}$ must land on some state twice.
    Say this happens for the first time at state $L$.
    \begin{center}
        \frame{\includegraphics[width=0.5\textwidth]{lecture-03/pumping-lemma-proof}}
    \end{center}
    Then taking $(x, y, z) := (S \to L, L \to L, L \to T)$ just works.
    ~ $\blacksquare$
\end{proof}

\begin{remark}
    All finite languages are regular.
    And this is trivially consistent with the pumping lemma---why?
\end{remark}

\begin{problem}
    Prove that the language $D := \{\str{0}^k \str{1}^k \mid k \geq 0\}$ is not regular.
\end{problem}

\begin{solution}
    Argue by contradiction.
    If the pumping length were $p$,
    then $s = \str{0}^p\str{1}^p \in D$ would contradict the pumping lemma.
    ~ $\blacksquare$
\end{solution}

\begin{problem}
    Prove that the language $B := \{w \mid w
    \text{ has as many } \str{0}\text{s} \text{ as } \str{1}\text{s}\}$
    is not regular.
\end{problem}

\begin{solution}
    The point is that if $B$ is regular,
    then so is $D = B \cap \str{0}^*\str{1}^*$ (by closure).
    Contradiction by the previous problem.
    ~ $\blacksquare$
\end{solution}

Unfortunately, there exist nonregular languages
that satisfy the pumping lemma;
take $\str{1}\str{1}^*D \cup \str{0}^*\str{1}^*$ with pumping length $1$, say.

\end{document}